\section{The guest instruction set}
\label{sec:app-isa}

This appendix lists the MIPS32 instructions a guest may execute, which is the
set the relation of \S\ref{sec:statement} quantifies over. It is also the
table from which the conformance corpus of \S\ref{sec:verification} is
generated: each row produces ten programs with pseudo-random operands,
boundary values and the delay-slot placements the instruction admits, giving
the 770 vectors compared against an independent emulator. The last column
names the chip that arithmetises the instruction, so the table fixes which
rows of Appendix~\ref{sec:app-chips} a program can occupy, and with it the
chip set that Section~\ref{sec:recursion} enumerates over.

Three conventions apply. \texttt{LO} and \texttt{HI} occupy the
register file as registers 32 and 33, so the four moves between them and the
general registers are additions of an immediate zero rather than a chip of
their own. \texttt{SYNC}, \texttt{SYNCI} and \texttt{PREF} are accepted and
retire as no-ops on \texttt{AddSubImm}, as an addition of zero. \texttt{SYSCALL} is absent
from the table because it is the boundary of the machine rather than a
computation within it; it is arithmetised by \texttt{SyscallInstrs} and described in
\S\ref{sec:ext-syscall}. MIPS32r2 user-mode integer
instructions outside the table (the conditional traps other than \texttt{teq},
the linking conditional branches \texttt{bgezal} and \texttt{bltzal}, the
branch-likely forms, \texttt{break} and \texttt{rdhwr}) are not supported: the
decoder maps them to an unimplemented opcode and execution stops with an error.

\paragraph{System calls the machine does not implement.} The behaviour is fixed in
every case; none is left unconstrained. A call number outside the machine's table
(neither an accelerator call nor a listed Linux number) stops execution with an error
before any state changes: the relation rejects, and no proof exists. Seven Linux calls
have their semantics (\texttt{brk}, \texttt{read}, \texttt{write},
\texttt{exit\_group}, \texttt{mmap}/\texttt{mmap2}, \texttt{clone},
\texttt{fcntl}). Seventeen more, the ones a standard runtime issues while starting
(\texttt{open}, \texttt{openat}, \texttt{close}, \texttt{fstat64},
\texttt{clock\_gettime}, \texttt{nanosleep}, \texttt{futex}, \texttt{prctl},
\texttt{uname}, \texttt{prlimit64}, \texttt{sigaltstack},
\texttt{rt\_sigaction}, \texttt{rt\_sigprocmask}, \texttt{madvise},
\texttt{munmap}, \texttt{gettid}, \texttt{sched\_getaffinity}), fall in two
groups, and the \texttt{SysLinux} table pins each outcome. Four whose result a program
would consume (\texttt{open}, \texttt{openat}, \texttt{fstat64},
\texttt{clock\_gettime}) fail as the kernel reports an unimplemented call:
\texttt{ENOSYS} in \texttt{\$v0} with the error flag \texttt{\$a3} set, no memory
touched, so the program sees the failure through its C library (a Rust program
reading the clock stops; a file open returns the error). The rest, issued by the
runtimes' start-up and ignored by them (Go's runtime stops when a signal call
fails), succeed with $0$ and touch nothing; none of them carries a value. The
accelerator calls are listed in Appendix~\ref{sec:app-chips}.

\footnotesize
\begin{longtable}{@{}llll>{\raggedright\arraybackslash}p{0.27\textwidth}l@{}}
\caption{The 77 guest instructions, their primary encoding fields and the chip that proves each. \emph{Op} is bits 31--26 and \emph{func} bits 5--0; a dash marks a field the format does not use, and the immediate forms are distinguished from the register forms by \emph{Op} alone. Rows that share both fields are told apart by another field: \texttt{bgez}, \texttt{bltz}, \texttt{bal} and \texttt{synci} by \emph{rt}, \texttt{rotr} and \texttt{rotrv} from \texttt{srl} and \texttt{srlv} by one bit of \emph{rs} or \emph{sa}, and \texttt{wsbh}, \texttt{seb} and \texttt{seh} by \emph{sa}. Semantics are as transcribed from the architecture manual into the specification table that generates the conformance corpus; on rows marked $^\dagger$ the implementation fixes a choice (\S\ref{sec:isa-emulator}): \texttt{j} and \texttt{jal} take the region bits \texttt{(PC + 4)[31:28]} as zero, exact for text below $2^{28}$, and \texttt{sc} always succeeds, setting \texttt{rt = 1}. Bit ranges \texttt{x[h:l]} include both ends.}\label{tab:isa}\\
\toprule
Mnemonic & Op & func & Family & Semantics & Chip \\
\midrule
\endfirsthead
\multicolumn{6}{@{}l}{\itshape Table~\ref{tab:isa}, continued from the previous page}\\
\toprule
Mnemonic & Op & func & Family & Semantics & Chip \\
\midrule
\endhead
\midrule
\multicolumn{6}{r@{}}{\itshape continued on the next page}\\
\endfoot
\bottomrule
\endlastfoot
\texttt{add} & \texttt{000000} & \texttt{100000} & Integer arithmetic & \texttt{rd = rs + rt} & \texttt{AddSub} \\
\texttt{addu} & \texttt{000000} & \texttt{100001} &  & \texttt{rd = rs + rt} & \texttt{AddSub} \\
\texttt{addi} & \texttt{001000} & -- &  & \texttt{rt = rs + sext(imm)} & \texttt{AddSubImm} \\
\texttt{addiu} & \texttt{001001} & -- &  & \texttt{rt = rs + sext(imm)} & \texttt{AddSubImm} \\
\texttt{sub} & \texttt{000000} & \texttt{100010} &  & \texttt{rd = rs - rt} & \texttt{AddSub} \\
\texttt{subu} & \texttt{000000} & \texttt{100011} &  & \texttt{rd = rs - rt} & \texttt{AddSub} \\
\texttt{mul} & \texttt{011100} & \texttt{000010} &  & \texttt{rd = rs * rt} & \texttt{Mul} \\
\texttt{mult} & \texttt{000000} & \texttt{011000} &  & \texttt{(hi, lo) = rs * rt} & \texttt{Mul} \\
\texttt{multu} & \texttt{000000} & \texttt{011001} &  & \texttt{(hi, lo) = rs * rt} & \texttt{Mul} \\
\texttt{div} & \texttt{000000} & \texttt{011010} &  & \texttt{(hi, lo) = (rs\%rt, rs/rt), signed} & \texttt{DivRem} \\
\texttt{divu} & \texttt{000000} & \texttt{011011} &  & \texttt{(hi, lo) = (rs\%rt, rs/rt), unsigned} & \texttt{DivRem} \\
\texttt{madd} & \texttt{011100} & \texttt{000000} &  & \texttt{(hi, lo) = (hi,lo) + rs * rt (signed)} & \texttt{MiscInstrs} \\
\texttt{maddu} & \texttt{011100} & \texttt{000001} &  & \texttt{(hi, lo) = (hi,lo) + rs * rt (unsigned)} & \texttt{MiscInstrs} \\
\texttt{msub} & \texttt{011100} & \texttt{000100} &  & \texttt{(hi, lo) = (hi,lo) - rs * rt (signed)} & \texttt{MiscInstrs} \\
\texttt{msubu} & \texttt{011100} & \texttt{000101} &  & \texttt{(hi, lo) = (hi,lo) - rs * rt (unsigned)} & \texttt{MiscInstrs} \\
\texttt{clo} & \texttt{011100} & \texttt{100001} &  & \texttt{rd = count\_leading\_ones(rs)} & \texttt{CloClz} \\
\texttt{clz} & \texttt{011100} & \texttt{100000} &  & \texttt{rd = count\_leading\_zeros(rs)} & \texttt{CloClz} \\
\addlinespace
\texttt{slt} & \texttt{000000} & \texttt{101010} & Comparison and logic & \texttt{rd = rs < rt} & \texttt{Lt} \\
\texttt{sltu} & \texttt{000000} & \texttt{101011} &  & \texttt{rd = rs < rt (unsigned)} & \texttt{Lt} \\
\texttt{slti} & \texttt{001010} & -- &  & \texttt{rt = rs < sext(imm)} & \texttt{LtImm} \\
\texttt{sltiu} & \texttt{001011} & -- &  & \texttt{rt = rs < sext(imm) (unsigned)} & \texttt{LtImm} \\
\texttt{and} & \texttt{000000} & \texttt{100100} &  & \texttt{rd = rs \& rt} & \texttt{Bitwise} \\
\texttt{andi} & \texttt{001100} & -- &  & \texttt{rt = rs \& zext(imm)} & \texttt{BitwiseImm} \\
\texttt{or} & \texttt{000000} & \texttt{100101} &  & \texttt{rd = rs | rt} & \texttt{Bitwise} \\
\texttt{ori} & \texttt{001101} & -- &  & \texttt{rt = rs | zext(imm)} & \texttt{BitwiseImm} \\
\texttt{xor} & \texttt{000000} & \texttt{100110} &  & \texttt{rd = rs \textasciicircum{} rt} & \texttt{Bitwise} \\
\texttt{xori} & \texttt{001110} & -- &  & \texttt{rt = rs \textasciicircum{} zext(imm)} & \texttt{BitwiseImm} \\
\texttt{nor} & \texttt{000000} & \texttt{100111} &  & \texttt{rd = \textasciitilde{}(rs | rt)} & \texttt{Bitwise} \\
\texttt{lui} & \texttt{001111} & -- &  & \texttt{rt = imm<{}<16} & \texttt{AddSubImm} \\
\addlinespace
\texttt{sll} & \texttt{000000} & \texttt{000000} & Shifts and bit fields & \texttt{rd = rt<{}<sa} & \texttt{ShiftLeftImm} \\
\texttt{sllv} & \texttt{000000} & \texttt{000100} &  & \texttt{rd = rt <{}< rs[4:0]} & \texttt{ShiftLeft} \\
\texttt{srl} & \texttt{000000} & \texttt{000010} &  & \texttt{rd = rt >{}> sa} & \texttt{ShiftRightImm} \\
\texttt{srlv} & \texttt{000000} & \texttt{000110} &  & \texttt{rd = rt >{}> rs[4:0]} & \texttt{ShiftRight} \\
\texttt{sra} & \texttt{000000} & \texttt{000011} &  & \texttt{rd = rt >{}> sa (arithmetic)} & \texttt{ShiftRightImm} \\
\texttt{srav} & \texttt{000000} & \texttt{000111} &  & \texttt{rd = rt >{}> rs[4:0] (arithmetic)} & \texttt{ShiftRight} \\
\texttt{rotr} & \texttt{000000} & \texttt{000010} &  & \texttt{rd = rotate\_right(rt, sa)} & \texttt{ShiftRightImm} \\
\texttt{rotrv} & \texttt{000000} & \texttt{000110} &  & \texttt{rd = rotate\_right(rt, rs[4:0])} & \texttt{ShiftRight} \\
\texttt{ext} & \texttt{011111} & \texttt{000000} &  & \texttt{rt = rs[msbd+lsb:lsb]} & \texttt{MiscInstrs} \\
\texttt{ins} & \texttt{011111} & \texttt{000100} &  & \texttt{rt = rt[31:msb+1] || rs[msb-lsb:0] || rt[lsb-1:0]} & \texttt{MiscInstrs} \\
\texttt{wsbh} & \texttt{011111} & \texttt{100000} &  & \texttt{rd = swaphalf(rt)} & \texttt{MovCond} \\
\texttt{seb} & \texttt{011111} & \texttt{100000} &  & \texttt{rd = signExtend(rt[7:0])} & \texttt{MiscInstrs} \\
\texttt{seh} & \texttt{011111} & \texttt{100000} &  & \texttt{rd = signExtend(rt[15:0])} & \texttt{MiscInstrs} \\
\addlinespace
\texttt{beq} & \texttt{000100} & -- & Control transfer & \texttt{PC = PC + 4 + sext(offset<{}<2), if rs == rt} & \texttt{Branch} \\
\texttt{bne} & \texttt{000101} & -- &  & \texttt{PC = PC + 4 + sext(offset<{}<2), if rs != rt} & \texttt{Branch} \\
\texttt{bgez} & \texttt{000001} & -- &  & \texttt{PC = PC + 4 + sext(offset<{}<2), if rs >= 0} & \texttt{Branch} \\
\texttt{bgtz} & \texttt{000111} & -- &  & \texttt{PC = PC + 4 + sext(offset<{}<2), if rs > 0} & \texttt{Branch} \\
\texttt{blez} & \texttt{000110} & -- &  & \texttt{PC = PC + 4 + sext(offset<{}<2), if rs <= 0} & \texttt{Branch} \\
\texttt{bltz} & \texttt{000001} & -- &  & \texttt{PC = PC + 4 + sext(offset<{}<2), if rs < 0} & \texttt{Branch} \\
\texttt{j} & \texttt{000010} & -- &  & \texttt{PC = (PC + 4)[31:28] || instr\_index || 00}$^\dagger$ & \texttt{Jump} \\
\texttt{jal} & \texttt{000011} & -- &  & \texttt{r31 = PC + 8, PC = (PC + 4)[31:28] || instr\_index || 00}$^\dagger$ & \texttt{Jump} \\
\texttt{jr} & \texttt{000000} & \texttt{001000} &  & \texttt{PC = rs} & \texttt{Jump} \\
\texttt{jalr} & \texttt{000000} & \texttt{001001} &  & \texttt{rd = PC + 8, PC = rs} & \texttt{Jump} \\
\texttt{bal} & \texttt{000001} & -- &  & \texttt{r31 = PC + 8, PC = PC + 4 + sext(offset<{}<2)} & \texttt{Jump} \\
\texttt{teq} & \texttt{000000} & \texttt{110100} &  & \texttt{trap, if rs == rt} & \texttt{MiscInstrs} \\
\addlinespace
\texttt{lb} & \texttt{100000} & -- & Aligned memory & \texttt{rt = sext(mem\_byte(base + offset))} & \texttt{LoadNarrow} \\
\texttt{lbu} & \texttt{100100} & -- &  & \texttt{rt = zext(mem\_byte(base + offset))} & \texttt{LoadNarrow} \\
\texttt{lh} & \texttt{100001} & -- &  & \texttt{rt = sext(mem\_halfword(base + offset))} & \texttt{LoadNarrow} \\
\texttt{lhu} & \texttt{100101} & -- &  & \texttt{rt = zext(mem\_halfword(base + offset))} & \texttt{LoadNarrow} \\
\texttt{lw} & \texttt{100011} & -- &  & \texttt{rt = mem\_word(base + offset)} & \texttt{LoadWord} \\
\texttt{ll} & \texttt{110000} & -- &  & \texttt{rt = mem\_word(base + offset)} & \texttt{LoadWord} \\
\texttt{sb} & \texttt{101000} & -- &  & \texttt{mem\_byte(base + offset) = rt} & \texttt{StoreNarrow} \\
\texttt{sh} & \texttt{101001} & -- &  & \texttt{mem\_halfword(base + offset) = rt} & \texttt{StoreNarrow} \\
\texttt{sw} & \texttt{101011} & -- &  & \texttt{mem\_word(base + offset) = rt} & \texttt{StoreWord} \\
\texttt{sc} & \texttt{111000} & -- &  & \texttt{mem\_word(base + offset) = rt, rt = 1, if atomic update, else rt = 0}$^\dagger$ & \texttt{StoreWord} \\
\addlinespace
\texttt{lwl} & \texttt{100010} & -- & Unaligned memory & \texttt{rt = rt merge most significant part of mem(base+offset)} & \texttt{MemoryUnaligned} \\
\texttt{lwr} & \texttt{100110} & -- &  & \texttt{rt = rt merge least significant part of mem(base+offset)} & \texttt{MemoryUnaligned} \\
\texttt{swl} & \texttt{101010} & -- &  & \texttt{store most significant part of rt} & \texttt{MemoryUnaligned} \\
\texttt{swr} & \texttt{101110} & -- &  & \texttt{store least significant part of rt} & \texttt{MemoryUnaligned} \\
\addlinespace
\texttt{mfhi} & \texttt{000000} & \texttt{010000} & Moves and no-ops & \texttt{rd = hi} & \texttt{AddSubImm} \\
\texttt{mflo} & \texttt{000000} & \texttt{010010} &  & \texttt{rd = lo} & \texttt{AddSubImm} \\
\texttt{mthi} & \texttt{000000} & \texttt{010001} &  & \texttt{hi = rs} & \texttt{AddSubImm} \\
\texttt{mtlo} & \texttt{000000} & \texttt{010011} &  & \texttt{lo = rs} & \texttt{AddSubImm} \\
\texttt{movn} & \texttt{000000} & \texttt{001011} &  & \texttt{rd = rs, if rt != 0} & \texttt{MovCond} \\
\texttt{movz} & \texttt{000000} & \texttt{001010} &  & \texttt{rd = rs, if rt == 0} & \texttt{MovCond} \\
\texttt{sync} & \texttt{000000} & \texttt{001111} &  & \texttt{sync (nop)} & \texttt{AddSubImm} \\
\texttt{synci} & \texttt{000001} & -- &  & \texttt{sync (nop)} & \texttt{AddSubImm} \\
\texttt{pref} & \texttt{110011} & -- &  & \texttt{prefetch(nop)} & \texttt{AddSubImm} \\
\addlinespace
\end{longtable}
\normalsize
